\documentclass[10pt]{extarticle}
\usepackage[a4paper,margin=1.5cm]{geometry}
\usepackage[T1]{fontenc}
\usepackage{amsmath,amssymb}
\usepackage{graphicx}
\usepackage[export]{adjustbox}
\usepackage{multicol}
\usepackage{enumitem}
\usepackage{booktabs,array,tabularx}
\setlength{\columnsep}{0.7cm}
\setlength{\parindent}{0pt}
\setlength{\parskip}{3pt}
\pagestyle{empty}
\makeatletter
\newcommand{\qfigdef}[4]{\@namedef{qf@file@#1}{#2}\@namedef{qf@w@#1}{#3}\@namedef{qf@wide@#1}{#4}}
\newcommand{\qf@check}[1]{\@ifundefined{qf@file@#1}{\PackageError{score2}{Unknown figure id #1}{}}{}}
\newcommand{\qf@img}[4][]{\includegraphics[width=#2\dimexpr\@nameuse{qf@w@#4}pt\relax,max width=#3,max height=0.45\textheight,#1]{\@nameuse{qf@file@#4}}}
\newif\ifqf@full
\newif\ifqf@boxed
\newsavebox{\qf@box}
\newcommand{\qf@open}{\global\qf@fulltrue\par\global\qf@boxedtrue
  \setbox\qf@box\vbox\bgroup\hsize\textwidth\linewidth\textwidth}
\newcommand{\qf@emit}[1]{\par\egroup\global\qf@boxedfalse
  \ifdim\dimexpr\ht\qf@box+\dp\qf@box\relax<#1\textheight\noindent\usebox{\qf@box}\par\else\unvbox\qf@box\fi}
\newcommand{\qfullsplit}{\ifqf@boxed\qf@emit{0.85}\fi}
\newcommand{\qf@close}{\ifqf@boxed\qf@emit{0.45}\fi\global\qf@fullfalse}
\newcommand{\qfullbegin}{\end{multicols}\qf@open}
\newcommand{\qfullend}{\qf@close\begin{multicols}{2}}
\newcommand{\qfullstart}{\qf@open}
\newcommand{\qfulldone}{\qf@close}
\newcommand{\qfig}[2][1]{\qf@check{#2}%
  \ifqf@full\qf@blk{#1}{#2}\else
  \ifnum\@nameuse{qf@wide@#2}=1 \expandafter\@firstoftwo\else\expandafter\@secondoftwo\fi
  {\end{multicols}\begin{center}\qf@img{#1}{\textwidth}{#2}\end{center}\begin{multicols}{2}}%
  {\qf@blk{#1}{#2}}\fi}
\newcommand{\qf@blk}[2]{\par\smallskip{\centering\qf@img{#1}{\linewidth}{#2}\par}\smallskip}
\newcommand{\qfigi}[2][1]{\qf@check{#2}\qf@img[valign=c]{#1}{\linewidth}{#2}}
\makeatother
\newcommand{\papersec}[1]{\par\bigskip{\centering\Large\bfseries PAPER-#1\par}\medskip}
\newcommand{\partsec}[2]{\par\medskip\noindent{\large\bfseries PART-#1: #2}\par\smallskip}
\newcommand{\sect}[1]{\par\smallskip\noindent\textbf{#1}\par}
\newcommand{\qn}[1]{\par\smallskip\noindent\textbf{#1.}~}
\newcommand{\paraq}[2]{\par\smallskip\noindent\textbf{Paragraph for Questions #1 and #2}\par}
\newenvironment{qoptions}{\begin{itemize}[leftmargin=1.8em,itemsep=1pt,topsep=2pt,label={}]}{\end{itemize}}
\newcommand{\opt}[2]{\item[(#1)] #2}
% figure registry (generated by score2_tools.py finalize)
\qfigdef{p1_phy_s1_q03_1}{Score-II - Test 07 (10 Jul 2022)-img/p1_phy_s1_q03_1.png}{56.3}{0}
\qfigdef{p1_phy_s1_q05_1}{Score-II - Test 07 (10 Jul 2022)-img/p1_phy_s1_q05_1.png}{72.0}{0}
\qfigdef{p1_phy_s1_q08_A}{Score-II - Test 07 (10 Jul 2022)-img/p1_phy_s1_q08_A.png}{111.1}{0}
\qfigdef{p1_phy_s1_q13_1}{Score-II - Test 07 (10 Jul 2022)-img/p1_phy_s1_q13_1.png}{86.7}{0}
\qfigdef{p1_phy_s1_q14_1}{Score-II - Test 07 (10 Jul 2022)-img/p1_phy_s1_q14_1.png}{111.9}{0}
\qfigdef{p1_phy_s1_q15_1}{Score-II - Test 07 (10 Jul 2022)-img/p1_phy_s1_q15_1.png}{123.2}{0}
\qfigdef{p1_phy_s1_q16_1}{Score-II - Test 07 (10 Jul 2022)-img/p1_phy_s1_q16_1.png}{151.4}{0}
\qfigdef{p1_phy_s1_q17_1}{Score-II - Test 07 (10 Jul 2022)-img/p1_phy_s1_q17_1.png}{90.0}{0}
\qfigdef{p1_phy_s1_q17_2}{Score-II - Test 07 (10 Jul 2022)-img/p1_phy_s1_q17_2.png}{89.4}{0}
\qfigdef{p1_phy_s1_q17_3}{Score-II - Test 07 (10 Jul 2022)-img/p1_phy_s1_q17_3.png}{88.3}{0}
\qfigdef{p1_phy_s1_q17_4}{Score-II - Test 07 (10 Jul 2022)-img/p1_phy_s1_q17_4.png}{89.2}{0}
\qfigdef{p1_phy_s1_q20_1}{Score-II - Test 07 (10 Jul 2022)-img/p1_phy_s1_q20_1.png}{93.0}{0}
\qfigdef{p1_phy_s2_q01_1}{Score-II - Test 07 (10 Jul 2022)-img/p1_phy_s2_q01_1.png}{111.5}{0}
\qfigdef{p1_phy_s2_q03_1}{Score-II - Test 07 (10 Jul 2022)-img/p1_phy_s2_q03_1.png}{95.5}{0}
\qfigdef{p1_phy_s2_q05_1}{Score-II - Test 07 (10 Jul 2022)-img/p1_phy_s2_q05_1.png}{89.2}{0}
\qfigdef{p1_phy_s2_q07_1}{Score-II - Test 07 (10 Jul 2022)-img/p1_phy_s2_q07_1.png}{46.3}{0}
\qfigdef{p1_phy_s2_q08_1}{Score-II - Test 07 (10 Jul 2022)-img/p1_phy_s2_q08_1.png}{110.6}{0}
\qfigdef{p1_phy_s2_q09_1}{Score-II - Test 07 (10 Jul 2022)-img/p1_phy_s2_q09_1.png}{139.8}{0}
\qfigdef{p1_chem_s1_q01_1}{Score-II - Test 07 (10 Jul 2022)-img/p1_chem_s1_q01_1.png}{103.9}{0}
\qfigdef{p1_chem_s1_q01_2}{Score-II - Test 07 (10 Jul 2022)-img/p1_chem_s1_q01_2.png}{51.1}{0}
\qfigdef{p1_chem_s1_q01_3}{Score-II - Test 07 (10 Jul 2022)-img/p1_chem_s1_q01_3.png}{60.7}{0}
\qfigdef{p1_chem_s1_q01_5}{Score-II - Test 07 (10 Jul 2022)-img/p1_chem_s1_q01_5.png}{32.5}{0}
\qfigdef{p1_chem_s1_q01_4}{Score-II - Test 07 (10 Jul 2022)-img/p1_chem_s1_q01_4.png}{57.8}{0}
\qfigdef{p1_chem_s1_q03_2}{Score-II - Test 07 (10 Jul 2022)-img/p1_chem_s1_q03_2.png}{53.7}{0}
\qfigdef{p1_chem_s1_q03_1}{Score-II - Test 07 (10 Jul 2022)-img/p1_chem_s1_q03_1.png}{39.2}{0}
\qfigdef{p1_chem_s1_q03_3}{Score-II - Test 07 (10 Jul 2022)-img/p1_chem_s1_q03_3.png}{46.8}{0}
\qfigdef{p1_chem_s1_q03_4}{Score-II - Test 07 (10 Jul 2022)-img/p1_chem_s1_q03_4.png}{46.6}{0}
\qfigdef{p1_chem_s1_q04_1}{Score-II - Test 07 (10 Jul 2022)-img/p1_chem_s1_q04_1.png}{89.2}{0}
\qfigdef{p1_chem_s1_q04_2}{Score-II - Test 07 (10 Jul 2022)-img/p1_chem_s1_q04_2.png}{90.6}{0}
\qfigdef{p1_chem_s1_q04_3}{Score-II - Test 07 (10 Jul 2022)-img/p1_chem_s1_q04_3.png}{89.7}{0}
\qfigdef{p1_chem_s1_q04_4}{Score-II - Test 07 (10 Jul 2022)-img/p1_chem_s1_q04_4.png}{90.8}{0}
\qfigdef{p1_chem_s1_q04_5}{Score-II - Test 07 (10 Jul 2022)-img/p1_chem_s1_q04_5.png}{90.9}{0}
\qfigdef{p1_chem_s1_q05_1}{Score-II - Test 07 (10 Jul 2022)-img/p1_chem_s1_q05_1.png}{30.6}{0}
\qfigdef{p1_chem_s1_q06_1}{Score-II - Test 07 (10 Jul 2022)-img/p1_chem_s1_q06_1.png}{79.6}{0}
\qfigdef{p1_chem_s1_q07_1}{Score-II - Test 07 (10 Jul 2022)-img/p1_chem_s1_q07_1.png}{71.8}{0}
\qfigdef{p1_chem_s1_q07_2}{Score-II - Test 07 (10 Jul 2022)-img/p1_chem_s1_q07_2.png}{54.2}{0}
\qfigdef{p1_chem_s1_q16_1}{Score-II - Test 07 (10 Jul 2022)-img/p1_chem_s1_q16_1.png}{105.6}{0}
\qfigdef{p1_chem_s2_q09_1}{Score-II - Test 07 (10 Jul 2022)-img/p1_chem_s2_q09_1.png}{135.6}{0}
\qfigdef{p1_math_s1_q12_1}{Score-II - Test 07 (10 Jul 2022)-img/p1_math_s1_q12_1.png}{137.4}{0}
\begin{document}
\begin{center}{\Large\bfseries JEE (Main+Advanced) Score-II --- Test 07 (10 Jul 2022)}\end{center}
\begin{multicols}{2}
\partsec{1}{PHYSICS}
\sect{Section-I}

\qn{1} A long cylinder with uniformly charged surface and cross sectional radius $a$ moves with constant velocity $v$ along its axis. Electric field strength at the surface of cylinder is equal to $E$. The current due to mechanical transfer of charge would be
\begin{qoptions}
\opt{A}{$EaV$}
\opt{B}{$2\pi\varepsilon_0 Eav$}
\opt{C}{$2\,Eav$}
\opt{D}{Zero}
\end{qoptions}

\qn{2} Eight small spheres of mercury, each with a diameter of 1 mm and a charge of 5 nC unite to form a large sphere. Potential energy stored in the large sphere is
\begin{qoptions}
\opt{A}{72 mJ}
\opt{B}{7.2 mJ}
\opt{C}{0.72 mJ}
\opt{D}{0.072 mJ}
\end{qoptions}

\qn{3} Two parallel plates A and B are having 12 $\mu$C and 4 $\mu$C charges. Their arrangement forms capacitance of 2 $\mu$F. It is connected with resistance of 1$\,\Omega$ at $t=0$. Choose CORRECT statement.
\qfig{p1_phy_s1_q03_1}
\begin{qoptions}
\opt{A}{Charge on the plate A is $4e^{-t/2}\ \mu$C, where $t$ is time in $\mu$S}
\opt{B}{Charge on the plate B is $-4e^{-t/2}\ \mu$C, where $t$ is time in $\mu$S}
\opt{C}{Charge on the plate A is $8+4e^{-t/2}\ \mu$C, where $t$ is time in $\mu$S}
\opt{D}{Charge on the plate B is $8+4e^{-t/2}\ \mu$C, where $t$ is time in $\mu$S}
\end{qoptions}

\qn{4} The dimensional formula for rate of change of electrical flux per unit current is
\begin{qoptions}
\opt{A}{$\left[\mathrm{ML^{3}T^{-4}A^{-2}}\right]$}
\opt{B}{$\left[\mathrm{ML^{2}T^{-4}A^{-2}}\right]$}
\opt{C}{$\left[\mathrm{ML^{3}T^{-2}A^{-2}}\right]$}
\opt{D}{$\left[\mathrm{ML^{3}T^{-2}A^{-1}}\right]$}
\end{qoptions}

\qn{5} Figure shows a regular hexagon made of uniform wire. If a battery is connected across the ends AB, magnetic field at O is $B_0$. If the same battery is connected across A and C, the magnetic field at O will be
\qfig{p1_phy_s1_q05_1}
\begin{qoptions}
\opt{A}{Equal to $B_0$}
\opt{B}{Greater than $B_0$}
\opt{C}{Less than $B_0$}
\opt{D}{May be less or greater than $B_0$}
\end{qoptions}

\qn{6} A paramagnetic material is kept in a magnetic field. The field is increased till the magnetization becomes constant. If the temperature is now decreased, the magnetization
\begin{qoptions}
\opt{A}{Will increase}
\opt{B}{Will decrease}
\opt{C}{Remains constant}
\opt{D}{May increase or decrease}
\end{qoptions}

\qn{7} A slab of thickness 10 cm is kept on $x$-$y$ plane ($z=0$) and its refractive index varies as $\mu=1.63\left[1+\dfrac{z}{20}\right]$ where $z$ is in cm. If it is viewed from top, its apparent thickness will be (in cm) $\left(\text{Given } \ln\dfrac{3}{2}=\dfrac{1.63}{4}\right)$
\begin{qoptions}
\opt{A}{5}
\opt{B}{6}
\opt{C}{4}
\opt{D}{3}
\end{qoptions}

\qn{8} Figure shows two coherent microwave sources $S_1$ and $S_2$ emitting waves of wavelength $\lambda$ and separated by a distance $3\lambda$. The minimum non-zero value of $y$ for point P to be an intensity maximum is ($D\gg\lambda$)
\qfig{p1_phy_s1_q08_A}
\begin{qoptions}
\opt{A}{$D$}
\opt{B}{$\sqrt{3}D$}
\opt{C}{$\dfrac{\sqrt{5}}{2}D$}
\opt{D}{$2\sqrt{2}D$}
\end{qoptions}

\qn{9} Suppose the price of a radioactive substance is directly proportional to the activity of radioactive sample. If initial price of radioactive sample (half life = 6 months) is Rs 4096, then price of radioactive sample after 5 years will be
\begin{qoptions}
\opt{A}{Rs 4}
\opt{B}{Rs 8}
\opt{C}{Rs 16}
\opt{D}{Rs 4096}
\end{qoptions}

\qn{10} A TV transmission tower at a particular station has a height 160 m. By how much the height of tower be increased to double its coverage range? Given radius of earth = 6400 km.
\begin{qoptions}
\opt{A}{320 m}
\opt{B}{160 m}
\opt{C}{480 m}
\opt{D}{640 m}
\end{qoptions}

\qn{11} A terrestrial telescope has an objective of focal length 180 cm and an eyepiece of focal length 5 cm. The erecting lens has a focal length of 3.5 cm. The separation between objective and eyepiece in normal adjustment is
\begin{qoptions}
\opt{A}{167 cm}
\opt{B}{179 cm}
\opt{C}{199 cm}
\opt{D}{213 cm}
\end{qoptions}

\qn{12} A proton moving with a speed $u$ directly towards an $\alpha$-particle (He nucleus) which, when the two particles were very far apart from one another, was initially at rest. What is the separation of the particles at their point of closest approach? (Assume $m_\alpha=4\,m_p$)
\begin{qoptions}
\opt{A}{$r=\dfrac{e^2}{\pi\varepsilon_0 m_p u^2}$}
\opt{B}{$r=\dfrac{e^2}{4\pi\varepsilon_0 m_p u^2}$}
\opt{C}{$r=\dfrac{5e^2}{4\pi\varepsilon_0 m_p u^2}$}
\opt{D}{$r=\dfrac{5e^2}{8\pi\varepsilon_0 m_p u^2}$}
\end{qoptions}

\qn{13} P-V diagram of cyclic process ABCA is as shown in figure. Choose the \textbf{CORRECT} statement(s):
\qfig{p1_phy_s1_q13_1}
\begin{qoptions}
\opt{A}{$\Delta Q_{AB}=+$ve}
\opt{B}{$\Delta U_{BC}=+$ve}
\opt{C}{$\Delta U_{CA}=-$ve}
\opt{D}{$\Delta W_{CAB}=+$ve}
\end{qoptions}

\qn{14} Two blocks A and B having equal mass, are placed in contact with each other on a rough plane, inclined at angle $\theta$ with horizontal as shown in figure. If coefficients of friction for A and B blocks are $\mu_1$ and $\mu_2$ ($\mu_1>\mu_2$) respectively, then for static equilibrium of blocks :-
\qfig{p1_phy_s1_q14_1}
\begin{qoptions}
\opt{A}{$\theta$ cannot be greater than $\tan^{-1}\mu_2$}
\opt{B}{$\theta$ cannot be less than $\tan^{-1}\mu_2$}
\opt{C}{maximum possible value of $\theta$ is equal to $\tan^{-1}\left(\dfrac{\mu_1+\mu_2}{2}\right)$}
\opt{D}{maximum possible value of $\theta$ is equal to $\tan^{-1}\sqrt{\mu_1\mu_2}$}
\end{qoptions}

\qn{15} In the siphon system in action shown in the figure $V$ refers to velocity, $P$ refers to pressure, then pick incorrect statement of the following (Consider ideal, non viscous fluid)
\qfig{p1_phy_s1_q15_1}
(Points A, B, C, D, E lie inside the pipe)\\
($P_0$ is atmospheric pressure)
\begin{qoptions}
\opt{A}{$V_A=V_B=V_E=\sqrt{2gh_1}$}
\opt{B}{$P_D=P_0-\rho g\,(h_1+h_2)$}
\opt{C}{$P_C>P_D$}
\opt{D}{$V_F=V_E$}
\end{qoptions}

\qn{16} In the graph shown, the potential energy of earth - satellite system is shown by a solid line as a function of distance $r$ (The separation between earth's centre and satellite). The total energy of the two objects which may or may not be bounded to the earth are shown in the figure by dotted lines.
\qfig{p1_phy_s1_q16_1}
Based on the above information, answer the following questions. Mark the correct statement(s).
\begin{qoptions}
\opt{A}{The object having the total energy $E_1$ is bounded one}
\opt{B}{The object having the total energy $E_2$ is bounded one}
\opt{C}{Both the objects are bounded}
\opt{D}{Both the objects are unbounded}
\end{qoptions}

\qn{17} Which one of the figure gives the temperature dependence of density of water correctly ?
\begin{qoptions}
\opt{A}{\qfigi{p1_phy_s1_q17_1}}
\opt{B}{\qfigi{p1_phy_s1_q17_2}}
\opt{C}{\qfigi{p1_phy_s1_q17_3}}
\opt{D}{\qfigi{p1_phy_s1_q17_4}}
\end{qoptions}

\qn{18} A listener is at rest with respect to the source of sound. A wind starts blowing along the line joining the source and the observer. Which of the following quantities do not change?
\begin{qoptions}
\opt{A}{Wavelength}
\opt{B}{Frequency}
\opt{C}{Angular wave number}
\opt{D}{Velocity of sound}
\end{qoptions}

\qn{19} A block is suspended by an ideal spring constant $K$. If the block is pulled down by constant force $F$ and if maximum displacement of block from it's initial position of rest is $z$, then
\begin{qoptions}
\opt{A}{$z=F/K$}
\opt{B}{$z=2F/K$}
\opt{C}{$z=3F/K$}
\opt{D}{$z=F/2K$}
\end{qoptions}

\qn{20} A solid sphere of mass $m$ and radius $R$ is lying on a rough horizontal plane. A constant force $4mg$ acts vertically at point P such that OP makes an angle $60^\circ$ with horizontal. Find the minimum value of coefficient of friction so that sphere starts pure rolling.
\qfig{p1_phy_s1_q20_1}
\begin{qoptions}
\opt{A}{$\dfrac{2}{7}$}
\opt{B}{$\dfrac{3}{7}$}
\opt{C}{$\dfrac{4}{7}$}
\opt{D}{$\dfrac{5}{7}$}
\end{qoptions}

\sect{Section-II}

\qn{1} An aluminium wire of length 60 cm is joined to a steel wire of length 80 cm and stretched between two fixed supports.
\qfig{p1_phy_s2_q01_1}
The tension produced is 2808 N. The cross-sectional areas of the steel and aluminium wires are 1.0 mm$^2$ and 3.0 mm$^2$ respectively. The densities of steel and aluminium are 70.2 g/cm$^3$ and 2.6 g/cm$^3$ respectively. Find the frequency (in Hz) of first overtone of this composite wire with the joint as a node

\qn{2} A parallel plate air capacitor is connected to a battery. The quantities-charge, potential difference, electric field and stored energy associated with the capacitor are $Q_0$, $V_0$, $E_0$ and $U_0$ respectively. A dielectric slab is now introduced with battery still connected. The dielectric constant of the dielectric is $k=2$. If final energy $U=pU_0$, find $p$.

\qn{3} Three particles A, B and C of equal mass move with equal speed $V=5\ \mathrm{ms^{-1}}$ along the medians of an equilateral triangle as shown in the figure. They collide at the centroid G of the triangle. After the collision, A comes to rest and B retraces its path with the speed $V=5\ \mathrm{ms^{-1}}$. What is the speed of C (in m/s) ?
\qfig{p1_phy_s2_q03_1}

\qn{4} The volume density of the elastic deformation energy in fresh water at a depth of $h=1$ km. (Take density of water $=10^3$ kg/m$^3$, $g=9.8$ m/s$^2$ and bulk modulus of elasticity of water $B=2\times10^{9}$ N/m$^2$) is $x\times10^{3}$ J/m$^3$. Find `$x$'.

\qn{5} Three capacitors each having capacitance $C=2\ \mu$F are connected with a battery of e.m.f. 30 V as shown in the figure. When the switch S is closed, find the amount of charge flown through the battery (in $\mu$C)
\qfig{p1_phy_s2_q05_1}

\qn{6} A block of mass 2 kg is attached to one end of a massless rod of length $\dfrac{1}{\pi}$ m. The rod is fixed to a horizontal plane at the other end such that the block and rod are free to revolve on a horizontal plane. The coefficient of friction between the block and surface is 0.1. Block is made to rotate with uniform speed by applying a constant external force in tangential direction on the block. Find the work done (in Joule) by external force when the rod rotates by $90^\circ$.

\qn{7} Three blocks A, B and C having masses $m$ kg, 2 kg and 3 kg are attached by massless strings and ideal pulleys as shown in the figure. The system is released from rest. If the block `A' remains stationary, find the value of $5m$ in kg.
\qfig{p1_phy_s2_q07_1}

\qn{8} A block of mass 1 kg is resting on a rough horizontal surface with coefficient of friction 0.2. If an impulse $\left(10\hat{i}-10\hat{j}\right)$ N-s is given to block then find velocity of block (in m/s) just after the given impulse
\qfig{p1_phy_s2_q08_1}

\qn{9} A block of mass 1 kg is pulled slowly slowly along the curve path ACB by a tangential force as shown in figure. Find the work done (in Joule) against the frictional force when the block moves from A to B.
\qfig{p1_phy_s2_q09_1}

\qn{10} Suppose a parent nucleus (which is being produced as the by-product of a nuclear reactor) decays into a daughter nucleus with a decay constant of $\lambda_p$ and the daughter nucleus decays with a decay constant of $\lambda_d$. At any instant the ratio of parent to daughter nuclei is observed to be 2. What should be the ratio $\dfrac{10\lambda_p}{\lambda_d}$ if the number of daughter nucleus is constant?

\partsec{2}{CHEMISTRY}
\sect{Section-I}

\qn{1} Which of the following product is not expected in given reaction ?
\qfig{p1_chem_s1_q01_1}
\begin{qoptions}
\opt{A}{\qfigi{p1_chem_s1_q01_2}}
\opt{B}{\qfigi{p1_chem_s1_q01_3}}
\opt{C}{\qfigi{p1_chem_s1_q01_5}}
\opt{D}{\qfigi{p1_chem_s1_q01_4}}
\end{qoptions}

\qn{2} Which reagent set should be used in given reaction sequence\\
Bromocyclopentane $\xrightarrow{\;?\;}\ \xrightarrow{\;?\;}$\\Trans-1,2-dibromocyclopentane
\begin{qoptions}
\opt{A}{Alco KOH/$\Delta$ and Br$_2$/H$_2$O}
\opt{B}{Conc. H$_2$SO$_4$/$\Delta$ and NBS/CCl$_4$}
\opt{C}{t-BuO--K$^{+}$/$\Delta$ and Br$_2$/H$_2$O}
\opt{D}{Alco. KOH/$\Delta$ and Br$_2$/CCl$_4$}
\end{qoptions}

\qn{3} Which of the following alcohol is expected to dehydrate using conc. H$_2$SO$_4$ at fastest rate ?
\begin{qoptions}
\opt{A}{\qfigi{p1_chem_s1_q03_2}}
\opt{B}{\qfigi{p1_chem_s1_q03_1}}
\opt{C}{\qfigi{p1_chem_s1_q03_3}}
\opt{D}{\qfigi{p1_chem_s1_q03_4}}
\end{qoptions}

\qn{4} What is the major organic product of the following reaction ?
\qfig{p1_chem_s1_q04_1}
\begin{qoptions}
\opt{A}{\qfigi{p1_chem_s1_q04_2}}
\opt{B}{\qfigi{p1_chem_s1_q04_3}}
\opt{C}{\qfigi{p1_chem_s1_q04_4}}
\opt{D}{\qfigi{p1_chem_s1_q04_5}}
\end{qoptions}

\qn{5} When a solution of isopropylamine is treated with NaNO$_2$ in HCl, which statement is INCORRECT?
\begin{qoptions}
\opt{A}{Formation of isopropyl alcohol will take place}
\opt{B}{Effervescence of N$_2$ gas will take place}
\opt{C}{Formation of \qfigi{p1_chem_s1_q05_1} takes place via carbocation formation}
\opt{D}{Formation of nitrosamines take place}
\end{qoptions}

\qn{6} Which of the following statement is correct ?
\begin{qoptions}
\opt{A}{S$_\mathrm{N}$2 reactivity of \qfigi{p1_chem_s1_q06_1}}
\opt{B}{Gutta-percha is cis-polyisoprene}
\opt{C}{Cellulose contains $\beta$-1,4 glycosidic linkages}
\opt{D}{CH$_3$OH and CH$_3$CH$_2$OH can not be distinguished using NaOH/I$_2$}
\end{qoptions}

\qn{7} Consider reaction sequence
\begin{center}\qfigi{p1_chem_s1_q07_1}\qfigi{p1_chem_s1_q07_2}\end{center}
Which of the following statement is INCORRECT ?
\begin{qoptions}
\opt{A}{Compound `A' can give Benzoic acid upon treatment with (i) KMnO$_4$/OH$^{-}$/$\Delta$ (ii) H$_3$O$^{\oplus}$}
\opt{B}{Compound `B' can give positive iodoform test}
\opt{C}{Compound `B' can give 3$^\circ$ alcohol upon treatment with (i) CH$_3$MgBr (ii) H$_3$O$^{\oplus}$}
\opt{D}{Compound `A' is more reactive in electrophilic aromatic substitution than toluene}
\end{qoptions}

\qn{8} Aqueous solution of CuSO$_4$ is treated with excess of aq. KCl. The formed compound and colour of solution is :
\begin{qoptions}
\opt{A}{CuS, black}
\opt{B}{$[\mathrm{CuCl_4}]^{3-}$, green-yellow}
\opt{C}{$[\mathrm{CuCl_4}]^{2-}$, bright green}
\opt{D}{Cu$_2$S, brown-black}
\end{qoptions}

\qn{9} For inducing the precipitation of Al(OH)$_3$ from a solution of sodium silicate and sodium aluminate, small amount of freshly prepared sample of which substance is added into the solution ?
\begin{qoptions}
\opt{A}{Na$_2$[Si(OH)$_4$]}
\opt{B}{NaOH}
\opt{C}{Hydrated Al$_2$O$_3$}
\opt{D}{SiC}
\end{qoptions}

\qn{10} The least stable halogen oxide among given is :
\begin{qoptions}
\opt{A}{BrO$_2$}
\opt{B}{ClO$_2$}
\opt{C}{I$_2$O$_4$}
\opt{D}{I$_2$O$_5$}
\end{qoptions}

\qn{11} Select the incorrect order of comparisons for the indicated properties :
\begin{qoptions}
\opt{A}{$\mathrm{Cu^{+2}(aq.)>Cu^{+}(aq.)}$ : stability}
\opt{B}{$\mathrm{Cr_2O_7^{2-}<MnO_4^{-}}$ : oxidising power}
\opt{C}{$\mathrm{Cr^{+2}>Fe^{+2}}$ : reducing power}
\opt{D}{C $>$ Fe $>$ Lanthanoid metal : percentage in mischmetal}
\end{qoptions}

\qn{12} The correct statement about halides of lead is :
\begin{qoptions}
\opt{A}{Lead (IV) chloride is highly stable towards heat.}
\opt{B}{Lead is known to form an iodide, PbI$_4$.}
\opt{C}{Lead (II) chloride reacts with Cl$_2$ to give PbCl$_4$.}
\opt{D}{Bonding in lead (IV) fluoride is typically covalent.}
\end{qoptions}

\qn{13} A sample of ground water has following contaminations :
\begin{center}
\begin{tabular}{|l|l|}\hline
Species & Concentration\\\hline
NO$_3^-$ & 35 ppm\\\hline
Pb & 20 ppb\\\hline
SO$_4^{2-}$ & 400 ppm\\\hline
Cd & 0.2 ppm\\\hline
\end{tabular}
\end{center}
This sample is unsuitable for drinking due to presence of :
\begin{qoptions}
\opt{A}{NO$_3^-$ ions}
\opt{B}{Pb}
\opt{C}{SO$_4^{2-}$ ions}
\opt{D}{Cd}
\end{qoptions}

\qn{14} 1 mole of XeF$_6$ reacts with sufficient amount of water to produce 80 g of HF along with a xenon compound `P'. Select the correct statement among given : (at. wt. of F = 19)
\begin{qoptions}
\opt{A}{The mentioned reaction depicts complete hydrolysis of XeF$_6$.}
\opt{B}{`P' has total 11 lone pair of electrons.}
\opt{C}{`P' is square pyramidal shaped.}
\opt{D}{The mentioned reaction is redox reaction.}
\end{qoptions}

\qn{15} 50 g of a solute is dissolved in 400 g of solvent (density = 1.0 g/ml) to make a solution with a volume of 420 ml. Select the correct statement about this solution.
\begin{qoptions}
\opt{A}{Molarity and molality of the solution are numerically equal.}
\opt{B}{Numerically, molarity is lower than molality.}
\opt{C}{Numerically, molarity is greater than molality.}
\opt{D}{Numerically, molarity and molality can't be compared without knowing the solute.}
\end{qoptions}

\qn{16} In the following graph, the natural log of the vapour pressures of two pure liquids (I and II) versus $\dfrac{1}{T}$ are plotted. What conclusion can be made ?
\qfig{p1_chem_s1_q16_1}
\begin{qoptions}
\opt{A}{$\Delta H_{\mathrm{vap.}}$ : liquid II $>$ liquid I}
\opt{B}{$\Delta H_{\mathrm{vap.}}$ : liquid II $=$ liquid I}
\opt{C}{$\Delta H_{\mathrm{vap.}}$ : liquid II $<$ liquid I}
\opt{D}{Enthalpies of vapourisation ($\Delta H_{\mathrm{vap.}}$) of these liquids can not be compared.}
\end{qoptions}

\qn{17} A metal (having molecular mass = 60) crystallises into a lattice containing a sequence of layer of ABC ABC close packing. If the distance between the two successive layer (A \& B) is 2.31 \AA. Calculate the density of unit cell.\\
($N_A=6\times10^{23}$) ($2.31\sqrt{3}=4.00$, $2.31\sqrt{2}=3.3$)
\begin{qoptions}
\opt{A}{6.25 kg/cm$^3$}
\opt{B}{0.625 g/cm$^3$}
\opt{C}{6.25 g/cm$^3$}
\opt{D}{6.25 g/m$^3$}
\end{qoptions}

\qn{18} Regarding criteria of catalysis which one of the following statement is not true ?
\begin{qoptions}
\opt{A}{The catalyst is unchanged chemically at the end of the chemical reaction.}
\opt{B}{A small quantity of catalyst is often sufficient to bring about a considerable amount of the reaction.}
\opt{C}{In reversible reaction, the catalyst alters the equilibrium position.}
\opt{D}{The catalyst accelerates the rate of reaction.}
\end{qoptions}

\qn{19} A solution is prepared by mixing 50.01 ml, 0.01 M, NaOH(aq.) and 50 ml, 0.005 M H$_2$SO$_4$(aq.). If few drops of phenolpthalein are now added into this solution, then find fraction (approx.) of this indicator present in its coloured form.\\
Given that :\\
Indicator constant of phenolpthalein is $2\times10^{-8}$\\
Ionic product of water is $10^{-14}$
\begin{qoptions}
\opt{A}{$\dfrac{1}{2}$}
\opt{B}{$\dfrac{1}{3}$}
\opt{C}{$\dfrac{2}{3}$}
\opt{D}{$\dfrac{10}{11}$}
\end{qoptions}

\qn{20} At what $\dfrac{[\mathrm{Br^-}]}{[\mathrm{Cl^-}]}$, does the following cell will be treated as dead cell ?\\
Ag(s)$|$AgCl(s)$|$NaCl(aq)$||$KBr(aq)$|$AgBr(s)$|$Ag(s)\\
Given that :\\
$K_{sp}$ of AgBr $=2\times10^{-13}$\\
$K_{sp}$ of AgCl $=8\times10^{-10}$
\begin{qoptions}
\opt{A}{$5\times10^{-3}$}
\opt{B}{$4\times10^{-4}$}
\opt{C}{$2\times10^{-3}$}
\opt{D}{$2.5\times10^{-4}$}
\end{qoptions}

\sect{Section-II}

\qn{1} How many structurally isomeric compounds can be formed using molecular formula C$_2$H$_2$O?

\qn{2} How many of the following amino acids are considered as basic amino acids ?\\
Aspartic acid, Serine, Tyrosine, Cysteine, Phenylalanine, Histidine, Valine, Arginine, Alanine

\qn{3} How many of the following are prepared by condensation polymerization ?\\
Dacron, Nylon-6, Nylon-66, PVC, SBR, PTFE, Bakelite, Novolac, Neoprene

\qn{4} Bond order of terminal S--O bond in pyrosulphate ion $\left(\mathrm{S_2O_7^{2-}}\right)$ is `x'. The value of 6`x' would be-

\qn{5} The number of optically active isomers for the complex species $[\mathrm{Co(gly)_2Cl_2}]^{-}$ are-

\qn{6} How many of the following mentioned minerals/substances are oxide or hydroxide ?\\
Cuprite, Calamine, Kaolinite, Bauxite, Magnetite, Zincite, Malachite, Cassiterite, Ruby gem stone, Copper matte, Siderite, Pig iron

\qn{7} A light source of wavelength $\lambda$ illuminates a metal and ejects photoelectron with $(\mathrm{K.E})_{\max}=1$ eV. Another light source of wavelength $\dfrac{\lambda}{3}$, ejects photoelectrons from same metal with $(\mathrm{K.E})_{\max}=5$ eV. Calculate the value of work function (eV) of metal.

\qn{8} Find $|\Delta_f G^\circ|$ (kcal/mole) of NH$_3$(g) at 300 K by using the following information :
\[2\mathrm{NH_3(g)}\ \xrightarrow[T\,=\,300\,\mathrm{K}]{\Delta_rH^\circ\,=\,40\ \mathrm{kcal/mole}}\ \mathrm{N_2(g)}+3\mathrm{H_2(g)}\]
\begin{center}
\begin{tabular}{|l|c|c|c|}\hline
 & N$_2$(g) & H$_2$(g) & NH$_3$(g)\\\hline
S$^\circ$(cal/K/mole) at 300 K & 50 & 30 & 40\\\hline
\end{tabular}
\end{center}

\qn{9} \qfig{p1_chem_s2_q09_1}
If molar ratio of CH$_4$(g) and He(g) in the shown container is 1 : 3 then the composition (mole ratio of He and CH$_4$) of gaseous mixture effusing out initially when valve is opened, is $x:1$. The value of $x$ is

\qn{10} For a reaction\\
$\mathrm{A(g)}\xrightarrow{\text{Rate constant}=k}\mathrm{P(g)}+\mathrm{Q(g)}$\\
carried in closed and rigid container at constant temperature $T$, following informations are provided.
\begin{center}
\begin{tabular}{|l|c|c|c|}\hline
 & Initial Concentration & $T_{1/2}$ & $T_{3/4}$\\\hline
A(g) & $1\times10^{-2}$ M & 10 Min & 30 Min\\\hline
\end{tabular}
\end{center}
If initial rate of reaction in mole/lit/min is represented in scientific notation as $x\times10^{-y}$ then find the value of $y$.

\partsec{3}{MATHEMATICS}
\sect{Section-I}

\qn{1} If the equation $x^3+ax^2+bx+216=0$ has three real roots in geometric progression, then $\dfrac{a^2+b^2}{ab}$ is equal to -
\begin{qoptions}
\opt{A}{$-\dfrac{37}{6}$}
\opt{B}{$-6$}
\opt{C}{$6$}
\opt{D}{$\dfrac{37}{6}$}
\end{qoptions}

\qn{2} Complete set of values of $m$ so that there is atleast one value of $x$ which satisfy the inequality $25^x+5^x<m$ is
\begin{qoptions}
\opt{A}{$m\in R$}
\opt{B}{$m\in\left(-\infty,-\dfrac{1}{4}\right)$}
\opt{C}{$m\in(0,\infty)$}
\opt{D}{$m\in\left(-\infty,\dfrac{1}{4}\right)$}
\end{qoptions}

\qn{3} A line makes an intercept of 5 and 4 on the $x$-axis \& $y$-axis respectively. If the axes are rotated about the origin by an angle $\alpha$ such that the length of intercept on the new $x$-axis is 8. The length of intercept on the new $y$-axis, is
\begin{qoptions}
\opt{A}{$\sqrt{23}$}
\opt{B}{$\dfrac{20}{8}$}
\opt{C}{$\dfrac{40}{\sqrt{139}}$}
\opt{D}{1}
\end{qoptions}

\qn{4} In the $\Delta$ABC, medians AD and BE are drawn. If AD $=4$, $\angle$BAD $=\pi/6$ and $\angle$ABE $=\pi/3$, then area of the $\Delta$ABC is
\begin{qoptions}
\opt{A}{$64/3$}
\opt{B}{$8/3$}
\opt{C}{$16/3$}
\opt{D}{$32/3\sqrt{3}$}
\end{qoptions}

\qn{5} Let $f(x)=x^4+ax^3+bx^2+cx+d$. If $ac-16d=0$, $g(x)=12x^5+15ax^4+20bx^3+30cx^2+60dx+k$, $a,b,c,d,k\in R$ has 5 positive real zeroes and $f(2)=1$, then the value of the sum of the zeroes of $f(x)$ if this sum is greater than 4 is
\begin{qoptions}
\opt{A}{12}
\opt{B}{15}
\opt{C}{20}
\opt{D}{30}
\end{qoptions}

\qn{6} Write the value of $\dfrac{\left(100\displaystyle\int_0^{\pi/2}\theta\sin^{10}\theta\,d\theta\right)-1}{\displaystyle\int_0^{\pi/2}\theta\sin^{8}\theta\,d\theta}$
\begin{qoptions}
\opt{A}{10}
\opt{B}{50}
\opt{C}{90}
\opt{D}{100}
\end{qoptions}

\qn{7} If $f(x)=4x+\pi\sin x$, $g(x)=f^{-1}(x)$ and $g''(3\pi)=\pi/n$, then write the value of `$|n|$'.
\begin{qoptions}
\opt{A}{4}
\opt{B}{16}
\opt{C}{32}
\opt{D}{64}
\end{qoptions}

\qn{8} If a complex number $z$ satisfies $|z+3|+|z-3|=10$, then the value of $\dfrac{60\,|z+3|}{\left|z+\bar z+(50/3)\right|}$ is
\begin{qoptions}
\opt{A}{18}
\opt{B}{30}
\opt{C}{36}
\opt{D}{60}
\end{qoptions}

\qn{9} The number of solutions of the equation $\sqrt2\cos3x-\cos2x=\sqrt3\sin2x+\sqrt2\sin3x$ in $\left[\dfrac{-4\pi}{3},\dfrac{2\pi}{3}\right]$ is
\begin{qoptions}
\opt{A}{3}
\opt{B}{4}
\opt{C}{5}
\opt{D}{6}
\end{qoptions}

\qn{10} If $y=\left(x-\sqrt{1+x^2}\right)^9$, then $(1+x^2)\dfrac{d^2y}{dx^2}+x\dfrac{dy}{dx}$ is
\begin{qoptions}
\opt{A}{$81y$}
\opt{B}{$-81y$}
\opt{C}{$162xy$}
\opt{D}{$-162xy$}
\end{qoptions}

\qn{11} In a binomial distribution B($n$, $p=3/4$) if the probability of at least one success is greater than or equal to $8/9$, then $n$ is greater than
\begin{qoptions}
\opt{A}{$\log_3 2$}
\opt{B}{$\log_2 5$}
\opt{C}{$\log_4 27$}
\opt{D}{$\log_2 3$}
\end{qoptions}

\qn{12} Length of the flagstaff, on the top of a vertical pole 4m high, subtends equal angles $30^\circ$ each at two points A, B on a straight line on the ground through the foot of the pole and on same side of pole. If distance of the point nearer to the foot of the pole is 4m, then length of the flagstaff is
\qfig{p1_math_s1_q12_1}
\begin{qoptions}
\opt{A}{$6\left(\sqrt3-1\right)$ m}
\opt{B}{$6\left(\sqrt3+1\right)$ m}
\opt{C}{$4\left(\sqrt3+1\right)$ m}
\opt{D}{None of these}
\end{qoptions}

\qn{13} Angle between the curves $y=\sqrt{|x|-1}$ and $y=\sqrt{5-x^2}$ at their points of intersection is -
\begin{qoptions}
\opt{A}{$\pi/6$}
\opt{B}{$\pi/3$}
\opt{C}{$\pi/2$}
\opt{D}{None of these}
\end{qoptions}

\qn{14} Solution set of $\cot^{-1}x>\sin^{-1}x$, is
\begin{qoptions}
\opt{A}{$-\infty<x<\sqrt{\dfrac{\sqrt5-2}{2}}$}
\opt{B}{$-1\leqslant x<\sqrt{\dfrac{\sqrt5-1}{2}}$}
\opt{C}{$-1\leqslant x<\sqrt{\dfrac{\sqrt5-2}{2}}$}
\opt{D}{None of these}
\end{qoptions}

\qn{15} Consider the statement p: There exists prime numbers that are not odd. Negation of statement `p' is
\begin{qoptions}
\opt{A}{There exist prime numbers that are odd}
\opt{B}{Not all prime numbers are even}
\opt{C}{All prime numbers are odd}
\opt{D}{None of these}
\end{qoptions}

\qn{16} A spherical iron-ball 20 cm in radius is coated with a layer of ice of uniform thickness that melts at a rate of 50 cm$^3$/min. When the thickness of ice is 5 cm, then the rate at which the thickness of ice decreases, is
\begin{qoptions}
\opt{A}{$\dfrac{1}{50\pi}$ cm/min}
\opt{B}{$\dfrac{1}{2\pi}$ cm/min}
\opt{C}{$\dfrac{1}{18\pi}$ cm/min}
\opt{D}{$\dfrac{1}{72\pi}$ cm/min}
\end{qoptions}

\qn{17} The set of points where $f(x)=\dfrac{\sin x}{1+|x|+x^2}$ is differentiable is
\begin{qoptions}
\opt{A}{$(-\infty,0)\cup(0,\infty)$}
\opt{B}{$(-\infty,-1)\cup(1,\infty)$}
\opt{C}{$(-\infty,-1)\cup(-1,\infty)$}
\opt{D}{$(-\infty,\infty)$}
\end{qoptions}

\qn{18} The differential equation of the family of circles passing through the origin and whose centre lies on the line $x=2$ is
\begin{qoptions}
\opt{A}{$(x^2-y^2+4x)\dfrac{dy}{dx}=2y(x-2)$}
\opt{B}{$(x^2+y^2+4x)\dfrac{dy}{dx}=2y(x-2)$}
\opt{C}{$(x^2-y^2-4x)\dfrac{dy}{dx}=2y(x-2)$}
\opt{D}{$(x^2+y^2-4x)\dfrac{dy}{dx}=2y(x-2)$}
\end{qoptions}

\qn{19} The probability of a shooter hitting a target is $2/3$. If he fires at a target 6 times, then the probability that he hits the target the second time in the fourth attempt is
\begin{qoptions}
\opt{A}{$8/27$}
\opt{B}{$16/81$}
\opt{C}{$4/27$}
\opt{D}{None of these}
\end{qoptions}

\qn{20} The value of $\sqrt2\displaystyle\int\dfrac{\cos x\,dx}{\sin\left(x-\frac{\pi}{4}\right)}$ is\\
(Where $c$ is arbitrary constant)
\begin{qoptions}
\opt{A}{$\log_e\left|\cos\left(x-\frac{\pi}{4}\right)\right|-x+c$}
\opt{B}{$\log_e\left|\cos\left(x-\frac{\pi}{4}\right)\right|+x+c$}
\opt{C}{$\log_e\left|\sin\left(x-\frac{\pi}{4}\right)\right|-x+c$}
\opt{D}{$\log_e\left|\sin\left(x-\frac{\pi}{4}\right)\right|+x+c$}
\end{qoptions}

\sect{Section-II}

\qn{1} For what value of $\lambda$, the system of equations $x+y+z=6$, $x+2y+3z=10$, $2x+4y+\lambda z=24$ is inconsistent :

\qn{2} Let $a=\log_3\log_3 4$. An integer $k$ satisfies $1<2^{\left(-k+3^{-a}\right)}<2$. Then $k$ equals

\qn{3} Let A and B be two square matrices of same order such that $A^3+B^3=0$ and $A^2B+AB^2=0$ but $|A|\neq|B|$. Then $|A^2-B^2|$ is equal to

\qn{4} The value of $\displaystyle\lim_{x\to0}\left(\dfrac{\displaystyle\int_0^{x^3}\sec^2t^2\,dt}{x\sin^2x}\right)$ is

\qn{5} Suppose $A_1,A_2,\dots,A_n$ are $n$ sets each having 5 elements and $B_1,B_2,\dots,B_{30}$ are thirty sets each with 3 elements, let $\displaystyle\bigcup_{i=1}^{n}A_i=\bigcup_{j=1}^{30}B_j=S$ and each element of $S$ belongs to exactly 10 of the $A_i$'s and exactly 9 of the $B_j$'s. Then $n$ is equal to

\qn{6} Let $A=\begin{pmatrix}1&-1&1\\2&1&-3\\1&1&1\end{pmatrix}$ and $10B=\begin{pmatrix}\alpha&2&2\\-5&0&\beta\\1&\gamma&3\end{pmatrix}$. If $B$ is the inverse of matrix $A$, then $\alpha$ is -

\qn{7} The sum of all odd divisors (odd natural numbers) of 2520 is

\qn{8} If $z^2+z+1=0$, where $z$ is complex number then the value of
\[\left(z+\frac1z\right)^3+\left(z^2+\frac1{z^2}\right)^3+\left(z^3+\frac1{z^3}\right)^3+\dots+\left(z^7+\frac1{z^7}\right)^3\]

\qn{9} If the area enclosed between the curve $y=\log_e(2x+e)$ and the coordinate axes is $K$, then $\dfrac1K$ is

\qn{10} If $\lambda\displaystyle\int_{-\pi/4}^{\pi/4}\dfrac{e^{\cos x}}{1+e^{\sin x}}\,dx=\int_{-\pi/4}^{\pi/4}e^{\cos x}\,dx$, then $\lambda$ equals -
\end{multicols}
\end{document}
